\section{Limitations and maturity}
\label{sec:limitations}

The honest summary of where this stands is: validated, not deployed. Concretely.

\paragraph{There is no accuracy claim in this document, anywhere.} No labelled
ground truth exists for the footage, so there is no intersection-over-union, no
precision or recall, no false-alarm rate and no detection-limit figure. Every
number in §\ref{sec:validation} is a throughput or a cost. A reader should treat
the absence as deliberate: the pipeline visibly separates a plume from a
background, and how well it does so against a human's judgement is unmeasured.

\paragraph{One camera, one scene type.} Everything here was developed against a
fixed thermal camera looking at one installation, and validated on synthetic
media built to resemble it. Generalisation to a different sensor, a different
spectral band, a moving camera or a cluttered outdoor scene is untested, and
several of the design choices --- a fixed region mask, a mean-of-$N$ background
--- assume a static camera and would need replacing rather than retuning.

\paragraph{A recipe is tuned to a scene, and does not travel.}
Figure~\ref{fig:plume}(e) is the demonstration: the shipped \texttt{motion}
recipe, run unmodified on a scene it was not tuned for, reports twenty-two
objects where a reader sees one plume. Nothing is broken --- the parameters are
simply wrong for that footage. This is the actual operating cost of the approach,
and the editor exists because of it: the tool makes retuning cheap, it does not
make it unnecessary.

\paragraph{Single-threaded, CPU only.} One core, one frame at a time. There is a
GPU in the benchmark machine and the engine never touches it. Frames are
processed in a single forward pass with no parallelism across frames, stages or
sinks, and no attempt has been made at any of the three.

\paragraph{Batch, not online.} The engine reads a file and writes files. There is
no streaming input from a live camera, no alerting, no persistence beyond CSV and
JSON on disk, and no operator-facing runtime --- the editor is a tuning tool, not
a control room display. Nothing in the architecture forbids these; none of them
is written.

\paragraph{No physical calibration.} A plume's segmented area is measured in
pixels, and its motion in pixels per frame. Converting either into a
concentration, a flow rate or a leak volume needs sensor calibration, scene
geometry and a physical model, none of which is in scope here.

\paragraph{Maturity.} By the TRL scale this is a level~4 result: the components
are validated together in a representative setting, on real footage, with a
reproducible measurement of what they cost. It is not hardened for on-site
deployment, has not run unattended, and has no failure-handling story beyond
raising an exception and stopping.
